% steric_sign_dcdt_walkthrough.tex
%
% Walk-through of how the steric chemical-potential term in the
% modified Nernst--Planck equation acts on dc_i/dt as the local
% concentration grows -- with the sign as written in
% docs/PNP Equation Formulations.tex (line 152), and the same
% calculation with the corrected Borukhov / Bazant sign.

\documentclass[11pt]{article}
\usepackage[margin=1in]{geometry}
\usepackage{amsmath,amssymb,bm}
\usepackage{xcolor}
\usepackage[most]{tcolorbox}
\usepackage{microtype}
\usepackage{parskip}
\usepackage{enumitem}
\usepackage{hyperref}
\usepackage{booktabs}

\setlist[itemize]{leftmargin=*,topsep=2pt,itemsep=3pt}
\setlist[enumerate]{leftmargin=*,topsep=2pt,itemsep=3pt}

\hypersetup{colorlinks=true, linkcolor=blue!50!black, urlcolor=blue!50!black}

\newcommand{\Deff}{D_{\mathrm{eff}}}
\newcommand{\steric}{\mathrm{steric}}

\title{What the steric term does to $\partial c_i/\partial t$\\
\large A walk-through of eq.~(15) of \texttt{PNP Equation Formulations.tex}}
\author{}
\date{2026-05-07}

\begin{document}
\maketitle

\section*{Setup}

The spec writes the modified Nernst--Planck equation
(\texttt{docs/PNP Equation Formulations.tex}, line 159) as
\begin{equation}
\label{eq:NPmod}
  \frac{\partial c_i}{\partial t}
  \;=\; \nabla\!\cdot\!\Bigl\{\, D_i \Bigl[\, \nabla c_i
        + c_i\, \nabla\!\Bigl(\tfrac{z_i F}{RT}\,\phi
                              + \mu^{\steric}(c)\Bigr)\Bigr]\Bigr\},
\end{equation}
with the steric chemical potential at line 152
\begin{equation}
\label{eq:muspec}
  \mu^{\steric}(c) \;=\; k_{B}T\, \ln\!\bigl(1 - \textstyle\sum_j a_j\, c_j\bigr).
\end{equation}

Two simplifications first:
\begin{itemize}
  \item Define the local packing fraction
        $\Phi \;\equiv\; \sum_j a_j\, c_j$.
  \item Switch to the dimensionless drift potential
        $\tilde\mu^{\steric} \equiv \mu^{\steric}/(k_{B}T) = \ln(1-\Phi)$.
        The units in \eqref{eq:muspec} are sloppy: what actually goes
        inside the gradient in \eqref{eq:NPmod} has to be dimensionless
        (it is added to $z_iF\phi/RT$), so it is $\tilde\mu^{\steric}$
        that enters, not $\mu^{\steric}$ as written. The qualitative
        story is identical.
\end{itemize}

The question we will answer is: \emph{as $c_i$ rises (and so $\Phi$
rises), what does the steric term do to $\partial c_i/\partial t$?}

\section{Step 1 -- write \eqref{eq:NPmod} as a drift--diffusion flux}

Writing \eqref{eq:NPmod} as $\partial c_i/\partial t = -\nabla\!\cdot J_i$
gives
\begin{align}
  J_i &= -D_i\Bigl[\nabla c_i
              + c_i\, \nabla\!\bigl(\tfrac{z_i F}{RT}\phi
                                   + \tilde\mu^{\steric}\bigr)\Bigr]
                                                       \nonumber \\[2pt]
      &= -D_i\, c_i\, \nabla\!\Bigl[\, \ln c_i
              + \tfrac{z_i F}{RT}\phi
              + \tilde\mu^{\steric}\,\Bigr].
  \label{eq:Jform}
\end{align}
The flux is $-D_i c_i$ times the gradient of a total drift potential.
The steric piece of that potential is just $\tilde\mu^{\steric} =
\ln(1-\Phi)$.

\section{Step 2 -- gradient of the steric potential}

Chain rule:
\begin{equation}
  \nabla \tilde\mu^{\steric}
   \;=\; \nabla\, \ln(1-\Phi)
   \;=\; \frac{1}{1-\Phi}\, \nabla(1-\Phi)
   \;=\; -\,\frac{\nabla\Phi}{1-\Phi}.
  \label{eq:gradmu}
\end{equation}
The steric \emph{contribution} to the flux of species $i$ is
\begin{equation}
  \Delta J_i^{\steric}
   \;=\; -D_i\, c_i\, \nabla\tilde\mu^{\steric}
   \;=\; +\, D_i\, c_i\, \frac{\nabla\Phi}{1-\Phi}.
  \label{eq:DeltaJ}
\end{equation}

Read the sign carefully. With the spec's $+\ln(1-\Phi)$, the steric flux
of species $i$ points \emph{up} the packing gradient -- toward whichever
side already has more counterion piled up. As $\Phi \to 1$ the prefactor
$1/(1-\Phi)$ diverges, so the up-gradient pump becomes infinitely
strong. This is already wrong before we put it back into
$\partial c_i/\partial t$.

\section{Step 3 -- collapse to a single-species effective diffusivity}

To see what \eqref{eq:NPmod} actually does to $\partial c_i/\partial t$
as $c_i$ grows, switch off electromigration for a moment and assume one
dominant counterion, so $\Phi = a_i\, c_i$ and $\nabla\Phi = a_i\, \nabla
c_i$. Combining the Fickian term and \eqref{eq:DeltaJ},
\begin{align}
  J_i &= -D_i\,\nabla c_i
        + D_i\, c_i\, \frac{a_i\, \nabla c_i}{1 - a_i c_i}
                                                       \nonumber \\[2pt]
      &= -D_i\,\nabla c_i\,
         \Bigl[\, 1 - \frac{a_i c_i}{1 - a_i c_i}\,\Bigr]
                                                       \nonumber \\[2pt]
      &= -D_i\,\nabla c_i\, \cdot \frac{1 - 2\Phi}{1 - \Phi}.
  \label{eq:Jspec}
\end{align}
So \eqref{eq:NPmod}, with the spec sign and electromigration off,
collapses to a nonlinear heat equation
\begin{equation}
\boxed{\;
  \frac{\partial c_i}{\partial t}
   \;=\; \nabla\!\cdot\!\bigl(\, \Deff(\Phi)\, \nabla c_i\,\bigr),
   \qquad
   \Deff(\Phi) \;=\; D_i\, \frac{1 - 2\Phi}{1 - \Phi}.
\;}
  \label{eq:DeffSpec}
\end{equation}
That $\Deff(\Phi)$ is the whole story.

\section{Step 4 -- what $\Deff$ does as $c_i$ grows}

Tabulate \eqref{eq:DeffSpec} versus packing fraction:

\begin{center}
\begin{tabular}{@{}lll@{}}
\toprule
$\Phi = a_i c_i$ & $\Deff(\Phi)$ & what eq.~(15) does \\
\midrule
$0$                  & $D_i$               & ordinary Fickian diffusion \\
$0 < \Phi < 1/2$     & $0 < \Deff < D_i$   & diffusion damped -- already a small artifact \\
$\Phi = 1/2$         & $0$                 & diffusion freezes -- no spatial relaxation \\
$1/2 < \Phi < 1$     & \textbf{negative}   & \textbf{backward heat equation} -- ill-posed \\
$\Phi \to 1$         & $-\infty$           & runaway: gradients amplified arbitrarily fast \\
\bottomrule
\end{tabular}
\end{center}

The pathology starts at $\Phi = 1/2$, \emph{not} at $\Phi = 1$. Anywhere
the counterion piles past half-packing, the spec's eq.~(15) is locally
a backward heat equation: the diffusive term \emph{amplifies} small
bumps in $c_i$ instead of smoothing them. There is no steady state for
it to find -- and that is exactly why Newton cannot solve the 4sp
warm-walk above $\psi_D \approx 4.83$ (nondim, $\approx +0.124$~V vs
RHE), which is where the ClO$_4^{-}$ EDL crosses $\Phi \approx 0.5$.
The solver is not failing for numerical reasons -- the PDE itself stops
being well-posed.

\section{Step 5 -- contrast with the corrected sign $-\ln(1-\Phi)$}

Repeat Steps 2--3 with $\tilde\mu^{\steric} = -\ln(1-\Phi)$
(Borukhov--Andelman--Orland 1997 eq.~(3); Bazant--Kilic--Storey--Ajdari
2009 eq.~(20)). Then $\nabla\tilde\mu^{\steric} = +\nabla\Phi/(1-\Phi)$
and
\begin{align}
  J_i &= -D_i\,\nabla c_i
        - D_i\, c_i\, \frac{a_i\,\nabla c_i}{1 - a_i c_i}
                                                       \nonumber \\[2pt]
      &= -D_i\,\nabla c_i\,
         \Bigl[\,1 + \frac{a_i c_i}{1 - a_i c_i}\,\Bigr]
                                                       \nonumber \\[2pt]
      &= -D_i\,\nabla c_i\, \cdot \frac{1}{1-\Phi}.
  \label{eq:Jcorrected}
\end{align}
Equivalently,
\begin{equation}
\boxed{\;
   \Deff(\Phi) \;=\; \frac{D_i}{1-\Phi}\;\;\;(\text{Borukhov sign}).
\;}
\end{equation}
Always positive, monotonically growing, and $\to +\infty$ as $\Phi \to
1$. Diffusion gets \emph{stronger} the closer the layer is to
saturation, which is precisely what enforces the Bikerman cap
$c_i \le 1/a_i$ (Borukhov 1997 eq.~(5)).

\section*{Summary}

\begin{itemize}
\item With the spec sign $+\ln(1-\Phi)$, the steric piece of the
      Nernst--Planck flux pumps mass \emph{up} the packing gradient.
      In the single-species reduction this gives
      $\Deff(\Phi) = D_i (1-2\Phi)/(1-\Phi)$, which is negative for
      $\Phi > 1/2$. Eq.~(15) then becomes a backward heat equation in
      the saturation skin -- ill-posed, no steady state, Newton
      diverges.
\item With the corrected Borukhov/Bazant sign $-\ln(1-\Phi)$, the
      steric piece of the flux pushes mass \emph{down} the packing
      gradient. The single-species reduction gives
      $\Deff(\Phi) = D_i/(1-\Phi)$, monotonically growing and
      diverging at saturation. The PDE is well-posed and produces the
      Bikerman cap $c_i \le 1/a_i$.
\end{itemize}

\paragraph{Caveats.} Two simplifications I papered over:
(i) the cross-species term in $\Phi = \sum_j a_j c_j$ produces
off-diagonal $c_i\,\nabla c_j/(1-\Phi)$ contributions that do not show
up in the single-species reduction, but they have the same sign
structure and do not rescue the spec form;
(ii) electromigration adds $c_i z_i F\,\nabla\phi/RT$ to the flux,
which does not enter $\Deff$. It only shifts \emph{where} $\Phi$ first
crosses $1/2$ -- it does not change that the equation goes ill-posed
once it does.

\paragraph{Where the fix lives in code.}
\begin{itemize}
  \item \texttt{Forward/bv\_solver/forms\_logc.py:303} ---
        \verb|mu_steric = -fd.ln(packing)|
  \item \texttt{Forward/bv\_solver/forms\_logc\_muh.py:347} --- same.
  \item Gated by \texttt{steric\_active} (line 283 / 335), which is
        \texttt{True} only when some $a_i \ne 0$ or a Bikerman
        counterion is configured. The 3sp+Boltzmann ideal preset
        ($a_{\mathrm{vals\_hat}} = [0,0,0]$, no
        \texttt{steric\_mode='bikerman'}) is unaffected.
\end{itemize}

The code is now consistent with the corrected sign; the spec doc at
\texttt{docs/PNP Equation Formulations.tex:152} still has the original
sign and disagrees with the code. The plan in
\texttt{docs/steric\_sign\_correction\_plan.md} flags this as needing
explicit advisor sign-off before the spec edit ships.

\end{document}
